% --- Archon physics formalization source begin ---
% archon:physics
% archon:covers PhyXMiniProblems/problem_phyx_mini_0657.lean
% archon:source-report reports/phyx_mini/problem_phyx_mini_0657.source.json

\chapter{Physics problem phyx\_mini\_0657}
\label{ch:PhyXMiniProblems_problem_phyx_mini_0657}

\paragraph{Problem source.}
\#\# Physical scenario

Figure shows the wave function of an electron in a rigid box. The electron energy is \$12.0\$ \$eV\$.

\#\# Question

What is the energy of the electron's ground state?

\#\# Answer choices

- A: A: 0.8 eV

- B: B: 1.1 eV

- C: C: 1.3 eV

- D: D: 1.7 eV

\#\# Figure caption (auxiliary; use the image as primary evidence)

The image is a graph depicting a wave function. The x-axis is labeled as \textbackslash{}( x \textbackslash{}), and the y-axis is labeled as \textbackslash{}( \textbackslash{}psi(x) \textbackslash{}). The graph shows a red sinusoidal wave that starts and ends at the y-axis baseline, with two peaks and one trough. This shape indicates a periodic function, suggesting it could represent a standing wave or similar phenomenon.

\#\# Dataset metadata

Quantum Phenomena; Physical Model Grounding Reasoning, Multi-Formula Reasoning

\paragraph{Recorded answer/context.}
C: C: 1.3 eV

\paragraph{Figure/image path.}
/root/proposal\_for\_physic/science-mango/phyx\_mini\_run/phyx\_data/test\_image/657.png

\paragraph{Formalization target.}
create a compiling Lean file with sorry bodies at `PhyXMiniProblems/problem\_phyx\_mini\_0657.lean`. The Lean declarations must preserve the physical quantities, dimensions or dimensional roles, figure labels, governing-law hypotheses, and final relation expressed by this problem.
Use Mathlib/Physlib names found through LeanExplore where available. If a physics API is missing, introduce faithful local abstractions rather than scalar placeholder aliases.

\begin{theorem}[Physics formalization target]
\label{thm:physics:phyx_mini_0657:target}
\lean{PhyXMiniProblems.ProblemPhyXMini0657.problem_phyx_mini_0657}
\uses{def:physics:phyx-mini-0657:phyxminiproblems-problemphyxmini0657-energyinelectronvolts, def:physics:phyx-mini-0657:phyxminiproblems-problemphyxmini0657-electronrigidboxsetup, def:physics:phyx-mini-0657:phyxminiproblems-problemphyxmini0657-matcheselectronrigidboxscenario, def:physics:phyx-mini-0657:phyxminiproblems-problemphyxmini0657-matchessuppliedwavefunctionfigure, def:physics:phyx-mini-0657:phyxminiproblems-problemphyxmini0657-matchesgivendepictedenergy, def:physics:phyx-mini-0657:phyxminiproblems-problemphyxmini0657-hasphysicalrigidboxparameters, def:physics:phyx-mini-0657:phyxminiproblems-problemphyxmini0657-satisfiesinfinitesquarewellnodallaw, def:physics:phyx-mini-0657:phyxminiproblems-problemphyxmini0657-satisfiesinfinitesquarewellenergyscaling, def:physics:phyx-mini-0657:phyxminiproblems-problemphyxmini0657-agreeswhenroundedtonearesttenthelectronvolt, def:physics:phyx-mini-0657:phyxminiproblems-problemphyxmini0657-isuniquenearestgroundstateenergychoice}
The assigned autoformalize agent should translate this physics problem into Lean declarations in the covered file, with theorem and lemma proof bodies written as `by sorry`.
\end{theorem}
\begin{proof}
This is an autoformalization task, not a proof task. Produce statements that can later be proved by the physics prover without weakening the physical model.
\end{proof}

% --- Archon declaration topology begin ---
\section{Declaration topology}
\begin{definition}[Length Quantity]
  \label{def:physics:phyx-mini-0657:phyxminiproblems-problemphyxmini0657-lengthquantity}
  \lean{PhyXMiniProblems.ProblemPhyXMini0657.LengthQuantity}
  A nonnegative, unit-independent physical length
\end{definition}
\begin{proof}
  This is the declared model definition of Length Quantity.
\end{proof}
\begin{definition}[length In Meters]
  \label{def:physics:phyx-mini-0657:phyxminiproblems-problemphyxmini0657-lengthinmeters}
  \lean{PhyXMiniProblems.ProblemPhyXMini0657.lengthInMeters}
  \uses{def:physics:phyx-mini-0657:phyxminiproblems-problemphyxmini0657-lengthquantity}
  Coherent-SI readout of a physical length, in metres
\end{definition}
\begin{proof}
  This is the declared model definition of length In Meters.
\end{proof}
\begin{definition}[energy In Joules]
  \label{def:physics:phyx-mini-0657:phyxminiproblems-problemphyxmini0657-energyinjoules}
  \lean{PhyXMiniProblems.ProblemPhyXMini0657.energyInJoules}
  Coherent-SI readout of a physical energy, in joules
\end{definition}
\begin{proof}
  This is the declared model definition of energy In Joules.
\end{proof}
\begin{definition}[energy In Electron Volts]
  \label{def:physics:phyx-mini-0657:phyxminiproblems-problemphyxmini0657-energyinelectronvolts}
  \lean{PhyXMiniProblems.ProblemPhyXMini0657.energyInElectronVolts}
  \uses{def:physics:phyx-mini-0657:phyxminiproblems-problemphyxmini0657-energyinjoules}
  Energy readout in electron volts, grounded by DimEnergy.electronVolt
\end{definition}
\begin{proof}
  This is the declared model definition of energy In Electron Volts.
\end{proof}
\begin{definition}[Particle Species]
  \label{def:physics:phyx-mini-0657:phyxminiproblems-problemphyxmini0657-particlespecies}
  \lean{PhyXMiniProblems.ProblemPhyXMini0657.ParticleSpecies}
  Particle species appearing in the problem statement
\end{definition}
\begin{proof}
  This is the declared model definition of Particle Species.
\end{proof}
\begin{definition}[Confinement Model]
  \label{def:physics:phyx-mini-0657:phyxminiproblems-problemphyxmini0657-confinementmodel}
  \lean{PhyXMiniProblems.ProblemPhyXMini0657.ConfinementModel}
  The idealized confinement model meant by a one-dimensional rigid box
\end{definition}
\begin{proof}
  This is the declared model definition of Confinement Model.
\end{proof}
\begin{definition}[Plot Axis]
  \label{def:physics:phyx-mini-0657:phyxminiproblems-problemphyxmini0657-plotaxis}
  \lean{PhyXMiniProblems.ProblemPhyXMini0657.PlotAxis}
  The two axes drawn in the supplied wave-function graph
\end{definition}
\begin{proof}
  This is the declared model definition of Plot Axis.
\end{proof}
\begin{definition}[Plot Axis Quantity]
  \label{def:physics:phyx-mini-0657:phyxminiproblems-problemphyxmini0657-plotaxisquantity}
  \lean{PhyXMiniProblems.ProblemPhyXMini0657.PlotAxisQuantity}
  Physical or mathematical role assigned to an axis label
\end{definition}
\begin{proof}
  This is the declared model definition of Plot Axis Quantity.
\end{proof}
\begin{definition}[Figure Node]
  \label{def:physics:phyx-mini-0657:phyxminiproblems-problemphyxmini0657-figurenode}
  \lean{PhyXMiniProblems.ProblemPhyXMini0657.FigureNode}
  The four zeroes visible from the left to the right box boundary
\end{definition}
\begin{proof}
  This is the declared model definition of Figure Node.
\end{proof}
\begin{definition}[Figure Lobe]
  \label{def:physics:phyx-mini-0657:phyxminiproblems-problemphyxmini0657-figurelobe}
  \lean{PhyXMiniProblems.ProblemPhyXMini0657.FigureLobe}
  The three consecutive lobes of the plotted stationary wave function
\end{definition}
\begin{proof}
  This is the declared model definition of Figure Lobe.
\end{proof}
\begin{definition}[Plot Curve Color]
  \label{def:physics:phyx-mini-0657:phyxminiproblems-problemphyxmini0657-plotcurvecolor}
  \lean{PhyXMiniProblems.ProblemPhyXMini0657.PlotCurveColor}
  Color of the curve in the supplied raster image
\end{definition}
\begin{proof}
  This is the declared model definition of Plot Curve Color.
\end{proof}
\begin{definition}[Rigid Box Wave Function Figure]
  \label{def:physics:phyx-mini-0657:phyxminiproblems-problemphyxmini0657-rigidboxwavefunctionfigure}
  \lean{PhyXMiniProblems.ProblemPhyXMini0657.RigidBoxWaveFunctionFigure}
  \uses{def:physics:phyx-mini-0657:phyxminiproblems-problemphyxmini0657-plotaxis, def:physics:phyx-mini-0657:phyxminiproblems-problemphyxmini0657-plotaxisquantity, def:physics:phyx-mini-0657:phyxminiproblems-problemphyxmini0657-figurenode, def:physics:phyx-mini-0657:phyxminiproblems-problemphyxmini0657-figurelobe, def:physics:phyx-mini-0657:phyxminiproblems-problemphyxmini0657-plotcurvecolor}
  Literal, uncalibrated graph data from image 657.png. Plot coordinates and amplitudes are real display coordinates, not substitutes for the physical box width or for the quantum state in the Hilbert space
\end{definition}
\begin{proof}
  This is the declared model definition of Rigid Box Wave Function Figure.
\end{proof}
\begin{definition}[Electron Rigid Box Setup]
  \label{def:physics:phyx-mini-0657:phyxminiproblems-problemphyxmini0657-electronrigidboxsetup}
  \lean{PhyXMiniProblems.ProblemPhyXMini0657.ElectronRigidBoxSetup}
  \uses{def:physics:phyx-mini-0657:phyxminiproblems-problemphyxmini0657-lengthquantity, def:physics:phyx-mini-0657:phyxminiproblems-problemphyxmini0657-particlespecies, def:physics:phyx-mini-0657:phyxminiproblems-problemphyxmini0657-confinementmodel, def:physics:phyx-mini-0657:phyxminiproblems-problemphyxmini0657-rigidboxwavefunctionfigure}
  Independent physical data for the depicted electron state. The state itself uses Physlib's one-dimensional quantum Hilbert space, while boxWidth and all level energies retain their physical dimensions
\end{definition}
\begin{proof}
  This is the declared model definition of Electron Rigid Box Setup.
\end{proof}
\begin{definition}[Matches Electron Rigid Box Scenario]
  \label{def:physics:phyx-mini-0657:phyxminiproblems-problemphyxmini0657-matcheselectronrigidboxscenario}
  \lean{PhyXMiniProblems.ProblemPhyXMini0657.MatchesElectronRigidBoxScenario}
  \uses{def:physics:phyx-mini-0657:phyxminiproblems-problemphyxmini0657-electronrigidboxsetup}
  The prose scenario: an electron stationary state in a one-dimensional rigid box
\end{definition}
\begin{proof}
  This is the declared model definition of Matches Electron Rigid Box Scenario.
\end{proof}
\begin{definition}[Matches Supplied Wave Function Figure]
  \label{def:physics:phyx-mini-0657:phyxminiproblems-problemphyxmini0657-matchessuppliedwavefunctionfigure}
  \lean{PhyXMiniProblems.ProblemPhyXMini0657.MatchesSuppliedWaveFunctionFigure}
  \uses{def:physics:phyx-mini-0657:phyxminiproblems-problemphyxmini0657-rigidboxwavefunctionfigure}
  Primary-image evidence: axes labelled x and psi(x), a red sinusoidal trace, four ordered boundary/node crossings, and the positive-negative-positive lobe pattern. This premise contains no energy level or answer-choice value
\end{definition}
\begin{proof}
  This is the declared model definition of Matches Supplied Wave Function Figure.
\end{proof}
\begin{definition}[Matches Given Depicted Energy]
  \label{def:physics:phyx-mini-0657:phyxminiproblems-problemphyxmini0657-matchesgivendepictedenergy}
  \lean{PhyXMiniProblems.ProblemPhyXMini0657.MatchesGivenDepictedEnergy}
  \uses{def:physics:phyx-mini-0657:phyxminiproblems-problemphyxmini0657-energyinelectronvolts, def:physics:phyx-mini-0657:phyxminiproblems-problemphyxmini0657-electronrigidboxsetup}
  The 12.0 eV energy belongs to the state depicted in the graph
\end{definition}
\begin{proof}
  This is the declared model definition of Matches Given Depicted Energy.
\end{proof}
\begin{definition}[Has Physical Rigid Box Parameters]
  \label{def:physics:phyx-mini-0657:phyxminiproblems-problemphyxmini0657-hasphysicalrigidboxparameters}
  \lean{PhyXMiniProblems.ProblemPhyXMini0657.HasPhysicalRigidBoxParameters}
  \uses{def:physics:phyx-mini-0657:phyxminiproblems-problemphyxmini0657-lengthinmeters, def:physics:phyx-mini-0657:phyxminiproblems-problemphyxmini0657-energyinjoules, def:physics:phyx-mini-0657:phyxminiproblems-problemphyxmini0657-electronrigidboxsetup}
  Positivity conditions selecting a nondegenerate physical rigid box
\end{definition}
\begin{proof}
  This is the declared model definition of Has Physical Rigid Box Parameters.
\end{proof}
\begin{definition}[Satisfies Infinite Square Well Nodal Law]
  \label{def:physics:phyx-mini-0657:phyxminiproblems-problemphyxmini0657-satisfiesinfinitesquarewellnodallaw}
  \lean{PhyXMiniProblems.ProblemPhyXMini0657.SatisfiesInfiniteSquareWellNodalLaw}
  \uses{def:physics:phyx-mini-0657:phyxminiproblems-problemphyxmini0657-electronrigidboxsetup}
  Nodal law for one-dimensional infinite-well stationary states: the nth state has n lobes and n - 1 interior nodes. It is a general governing relation between the raw figure counts and the state's independent quantum number, not an assumption that the depicted state is specifically n = 3
\end{definition}
\begin{proof}
  This is the declared model definition of Satisfies Infinite Square Well Nodal Law.
\end{proof}
\begin{definition}[Satisfies Infinite Square Well Energy Scaling]
  \label{def:physics:phyx-mini-0657:phyxminiproblems-problemphyxmini0657-satisfiesinfinitesquarewellenergyscaling}
  \lean{PhyXMiniProblems.ProblemPhyXMini0657.SatisfiesInfiniteSquareWellEnergyScaling}
  \uses{def:physics:phyx-mini-0657:phyxminiproblems-problemphyxmini0657-energyinelectronvolts, def:physics:phyx-mini-0657:phyxminiproblems-problemphyxmini0657-electronrigidboxsetup}
  Energy scaling for a particle in a fixed one-dimensional infinite square well: E\_n = n\textasciicircum{}2 E\_1 for each positive integer n. This is the smallest local governing-law interface needed here; it neither fixes E\_1 nor mentions the displayed answer
\end{definition}
\begin{proof}
  This is the declared model definition of Satisfies Infinite Square Well Energy Scaling.
\end{proof}
\begin{lemma}[supplied Figure depicts third Stationary State]
  \label{lem:physics:phyx-mini-0657:phyxminiproblems-problemphyxmini0657-suppliedfigure-depicts-thirdstationarystate}
  \lean{PhyXMiniProblems.ProblemPhyXMini0657.suppliedFigure_depicts_thirdStationaryState}
  \uses{def:physics:phyx-mini-0657:phyxminiproblems-problemphyxmini0657-electronrigidboxsetup, def:physics:phyx-mini-0657:phyxminiproblems-problemphyxmini0657-matchessuppliedwavefunctionfigure, def:physics:phyx-mini-0657:phyxminiproblems-problemphyxmini0657-satisfiesinfinitesquarewellnodallaw}
  The two interior nodes and three lobes identify the plotted state as n = 3
\end{lemma}
\begin{proof}
  The conclusion follows from the governing relations and figure readouts named in the statement of supplied Figure depicts third Stationary State.
\end{proof}
\begin{definition}[Answer Choice]
  \label{def:physics:phyx-mini-0657:phyxminiproblems-problemphyxmini0657-answerchoice}
  \lean{PhyXMiniProblems.ProblemPhyXMini0657.AnswerChoice}
  Labels of the four ground-state-energy choices printed in the problem
\end{definition}
\begin{proof}
  This is the declared model definition of Answer Choice.
\end{proof}
\begin{definition}[energy Electron Volts]
  \label{def:physics:phyx-mini-0657:phyxminiproblems-problemphyxmini0657-answerchoice-energyelectronvolts}
  \lean{PhyXMiniProblems.ProblemPhyXMini0657.AnswerChoice.energyElectronVolts}
  \uses{def:physics:phyx-mini-0657:phyxminiproblems-problemphyxmini0657-answerchoice}
  Energy printed beside each answer choice, in electron volts
\end{definition}
\begin{proof}
  This is the declared model definition of energy Electron Volts.
\end{proof}
\begin{definition}[recorded Dataset Answer]
  \label{def:physics:phyx-mini-0657:phyxminiproblems-problemphyxmini0657-recordeddatasetanswer}
  \lean{PhyXMiniProblems.ProblemPhyXMini0657.recordedDatasetAnswer}
  \uses{def:physics:phyx-mini-0657:phyxminiproblems-problemphyxmini0657-answerchoice}
  Dataset metadata recording answer C; deliberately not used as a premise
\end{definition}
\begin{proof}
  This is the declared model definition of recorded Dataset Answer.
\end{proof}
\begin{definition}[ground State Choice Error Electron Volts]
  \label{def:physics:phyx-mini-0657:phyxminiproblems-problemphyxmini0657-groundstatechoiceerrorelectronvolts}
  \lean{PhyXMiniProblems.ProblemPhyXMini0657.groundStateChoiceErrorElectronVolts}
  \uses{def:physics:phyx-mini-0657:phyxminiproblems-problemphyxmini0657-energyinelectronvolts, def:physics:phyx-mini-0657:phyxminiproblems-problemphyxmini0657-electronrigidboxsetup, def:physics:phyx-mini-0657:phyxminiproblems-problemphyxmini0657-answerchoice}
  Absolute error between the physical ground energy and a displayed choice
\end{definition}
\begin{proof}
  This is the declared model definition of ground State Choice Error Electron Volts.
\end{proof}
\begin{definition}[Agrees When Rounded To Nearest Tenth Electron Volt]
  \label{def:physics:phyx-mini-0657:phyxminiproblems-problemphyxmini0657-agreeswhenroundedtonearesttenthelectronvolt}
  \lean{PhyXMiniProblems.ProblemPhyXMini0657.AgreesWhenRoundedToNearestTenthElectronVolt}
  \uses{def:physics:phyx-mini-0657:phyxminiproblems-problemphyxmini0657-electronrigidboxsetup, def:physics:phyx-mini-0657:phyxminiproblems-problemphyxmini0657-answerchoice, def:physics:phyx-mini-0657:phyxminiproblems-problemphyxmini0657-groundstatechoiceerrorelectronvolts}
  Agreement with an energy displayed to the nearest tenth of an electron volt
\end{definition}
\begin{proof}
  This is the declared model definition of Agrees When Rounded To Nearest Tenth Electron Volt.
\end{proof}
\begin{definition}[Is Unique Nearest Ground State Energy Choice]
  \label{def:physics:phyx-mini-0657:phyxminiproblems-problemphyxmini0657-isuniquenearestgroundstateenergychoice}
  \lean{PhyXMiniProblems.ProblemPhyXMini0657.IsUniqueNearestGroundStateEnergyChoice}
  \uses{def:physics:phyx-mini-0657:phyxminiproblems-problemphyxmini0657-electronrigidboxsetup, def:physics:phyx-mini-0657:phyxminiproblems-problemphyxmini0657-answerchoice, def:physics:phyx-mini-0657:phyxminiproblems-problemphyxmini0657-groundstatechoiceerrorelectronvolts}
  The selected choice is strictly nearer than every other displayed energy
\end{definition}
\begin{proof}
  This is the declared model definition of Is Unique Nearest Ground State Energy Choice.
\end{proof}
% --- Archon declaration topology end ---
% --- Archon physics formalization source end ---
